\chapter{Будущие ветви: как расширять систему, не теряя её}

\section{Две обещанные дальние ветви}

У расширенного курса есть ещё одна задача: не просто завершить сегодняшнюю версию,
но и подготовить язык, на котором будут обсуждаться следующие большие ветви. На
горизонте уже видны по меньшей мере две:
\begin{enumerate}
\item резервуарные вычисления как дополнительная ассоциативно-динамическая среда;
\item ассоциативная память на кубитах как исследовательская ветвь квантового слоя.
\end{enumerate}

\section{Что нужно требовать от любой новой ветви}

Новая ветвь не должна приходить в систему как эстетический жест. Для неё нужен
одинаковый протокол допуска.
\begin{enumerate}
\item Какая конкретная слабость базовой архитектуры атакуется?
\item Как будет измеряться улучшение?
\item Как ветвь взаимодействует с provenance, deletion и revision?
\item Что остаётся в явном графе, а что уходит в скрытую динамику?
\end{enumerate}

\section{Резервуарная ветвь: что в ней действительно интересно}

Резервуарная схема интересна не тем, что ``она нейроподобна'', а тем, что даёт
богатый нелинейный feature space для частичных cue и сложных временных контекстов.
Если использовать её осторожно, она может:
\begin{enumerate}
\item улучшать pattern completion;
\item давать признаки для сложных mood/history-dependent retrieval policies;
\item служить дополнительным предложителем bridge-кандидатов.
\end{enumerate}

Но её нельзя пускать напрямую в write-path без жёсткого интерфейса.

\section{Квантовая ветвь: чего нельзя обещать заранее}

Квантовая ассоциативная память соблазнительна по двум причинам: волновая интуиция
и естественная работа с суперпозициями. Но сама по себе эта красота ничего не
доказывает. До стенда допустимы только такие утверждения:
\begin{enumerate}
\item ветвь может дать иной режим диффузии и поиска;
\item ветвь может дать иной bias к pattern completion;
\item ветвь \emph{не считается полезной}, пока не показала измеримый выигрыш на
общем стенде против классических альтернатив.
\end{enumerate}

\section{Правильная амбиция}

\begin{principle}[Большая ветвь должна быть сравнима с малой системой]
Чем радикальнее новая ветвь, тем строже она обязана сравниваться с уже работающей
графовой памятью. Иначе вместо расширения курса получается побег из проверяемости.
\end{principle}

Это правило и будет держать будущие главы честными.
